\SourcePage{246}{249}
\section{Spin}

In classical mechanics as well as in quantum mechanics, conservation of
angular momentum follows from the isotropy of space for a closed system. This
already reveals the connection between angular momentum and rotational
symmetry. In quantum mechanics, however, that connection becomes especially
fundamental: it essentially supplies the principal content of the notion of
angular momentum. This is all the more so because the classical definition of
a particle's angular momentum as the product
$\mathbf r\times\mathbf p$ loses its direct meaning here, since the position
vector and momentum cannot be measured simultaneously.

As we saw in \S~28, the values of $l$ and $m$ determine the angular
dependence of a particle's wave function and, with it, all the rotational
symmetry properties of that function. In their most general form these
properties are specified by stating how wave functions transform when the
coordinate system is rotated.

The wave function $\psi_{LM}$ of a system of particles, with prescribed
angular momentum $L$ and projection $M$, is unchanged\footnote{Apart from an
inessential phase factor.} only by rotations of the coordinate system about
the $z$ axis. After any rotation that alters the direction of this axis, the
projection of the angular momentum on $z$ no longer has a definite value. In
the new coordinate axes the wave function must therefore become, in general,
a superposition (linear combination) of the $2L+1$ functions associated with
the possible values of $M$ at fixed $L$. Equivalently, rotations of the
coordinate system transform the $2L+1$ functions $\psi_{LM}$ into one
another\footnote{In mathematical terminology, these functions realize an
\emph{irreducible representation} of the rotation group. The number of
functions transformed into one another is called the \emph{dimension of the
representation}; it is understood that no choice of other linear
combinations of these functions can reduce this number.}. The transformation
law---that is, the superposition coefficients as functions of the angles
through which the coordi\SourcePageJoin{247}{250}nate axes are rotated---is
fixed completely by the value of $L$. Angular momentum thus takes on the
meaning of a quantum number that classifies the states of a system by their
transformation properties under rotations of the coordinate system. This
aspect of angular momentum is especially important in quantum mechanics
because it does not depend directly on an explicit angular dependence of the
wave functions: their mutual transformation law can be formulated by itself,
without referring to such a dependence.


Consider a composite particle (an atomic nucleus, for example) that is at rest
as a whole and occupies a definite internal state. Besides having a definite
internal energy, it has angular momentum of definite magnitude $L$, arising
from the motion of the particles within it; this angular momentum can still
have $2L+1$ distinct orientations in space. In other words, when treating the
motion of the composite particle as a whole, we must assign to it, in addition
to its coordinates, one discrete variable: the projection of its internal
angular momentum on some selected direction in space.

With the meaning of angular momentum just described, however, its origin
ceases to be essential. We are therefore led naturally to the idea of an
``intrinsic'' angular momentum that must be assigned to a particle regardless
of whether it is ``composite'' or ``elementary.''

Thus quantum mechanics requires an elementary particle to possess a certain
``intrinsic'' angular momentum unrelated to its motion through space. This
property of elementary particles is specifically quantum mechanical (it
vanishes in the limit $\hbar\to0$) and admits no classical interpretation in
principle.\footnote{In particular, it would be entirely meaningless to picture
the ``intrinsic'' angular momentum of an elementary particle as resulting from
its rotation ``about its own axis.''}

The intrinsic angular momentum of a particle is called its \emph{spin}, in
contrast to the angular momentum associated with the particle's motion in
space, which is called \emph{orbital angular momentum}.\footnote{The physical
idea that the electron has an intrinsic angular momentum was proposed by
Uhlenbeck and Goudsmit (G.~Uhlenbeck, S.~Goudsmit, 1925). Spin was introduced
into quantum mechanics by Pauli (W.~Pauli, 1927).} The particle in question
may be elementary, or it may be composite but behave as an elementary object
within the class of phenomena under consideration (an atomic nucleus, for
example). We shall denote the particle's spin, measured like orbital angular
momentum in units of $\hbar$, by $s$.

\SourcePage{248}{251}
For a particle with spin, its wave function must specify not only the
probabilities of its possible positions in space but also the probabilities
of the possible orientations of its spin. In other words, the wave function
must depend not merely on three continuous variables---the particle's
coordinates---but also on one discrete \emph{spin variable}. This variable
gives the value of the spin projection on a selected spatial direction (the
$z$ axis) and runs through a finite set of discrete values, which we shall
subsequently denote by $\sigma$.

Let $\psi(x,y,z;\sigma)$ be such a wave function. In substance it is a set of
several different functions of the coordinates, one for each value of
$\sigma$; we shall call these functions the \emph{spin components} of the wave
function. The integral
\[
 \int |\psi(x,y,z;\sigma)|^2\,dV
\]
then gives the probability that the particle has the specified value of
$\sigma$. The probability of finding the particle in the volume element $dV$
with an arbitrary value of $\sigma$ is
\[
 dV\sum_\sigma |\psi(x,y,z;\sigma)|^2.
\]

When the quantum-mechanical spin operator is applied to a wave function, it
acts specifically on the spin variable $\sigma$. In other words, it effects
some transformation among the components of the wave function. Its explicit
form will be found below. Even without that form, however, very general
considerations suffice to show that $\hat s_x$, $\hat s_y$, and $\hat s_z$
obey the same commutation relations as the orbital-angular-momentum
operators.


An angular-momentum operator is essentially identical with the operator of an
infinitesimal rotation. In deriving the orbital-angular-momentum operator in
\S~26, we considered how a rotation acts on a function of the coordinates.
That derivation is inapplicable to spin angular momentum, since the spin
operator acts on the spin variable rather than on the coordinates. To obtain
the required commutation relations, we must consequently consider an
infinitesimal rotation in its general form, as a rotation of the coordinate
system. Applying infinitesimal rotations in succession, first about the $x$ and $y$
axes and then about the same axes in reverse order,

we readily find by direct computation that

the difference between the outcomes of these two operations is equivalent to an
infinitesimal \SourcePage{249}{252}rotation about the $z$ axis

(through an angle equal to the product of the rotation angles about the $x$ and
$y$ axes)
.
We shall omit these elementary calculations here. They lead once more to the
standard commutation relations for the component operators of angular momentum;
the same relations must therefore hold for the spin operators:

\begin{equation}
 \{\hat s_x,\hat s_y\}=i\hat s_z
 \tag{54.1}
\end{equation}
together with all their physical consequences.

The commutation relations (54.1) allow the possible values of the magnitude
and components of spin to be determined. The entire derivation in \S~27
(equations (27.7)--(27.9)) relied only on the commutation relations and is
therefore fully applicable here, with $s$ replacing $L$ in those equations.
It follows from equations (27.7) that the eigenvalues of a spin projection
form a sequence whose adjacent terms differ by unity. We can no longer assert,
however, that these values themselves must be integers, as was the case for
the orbital-angular-momentum projection $L_z$. The argument at the beginning
of \S~27 does not apply here, since it rests on the expression (26.14) for the
operator $\hat l_z$, which is specific to orbital angular momentum.

Furthermore, the sequence of eigenvalues of $s_z$ is bounded above and below
by values equal in magnitude and opposite in sign, which we denote by
$\pm s$. The difference $2s$ between the largest and smallest values of $s_z$
must be an integer or zero. Hence $s$ may take the values $0$, $1/2$, $1$,
$3/2$,~\dots

Thus the eigenvalues of the square of the spin are
\begin{equation}
 s^2=s(s+1),
 \tag{54.2}
\end{equation}
where $s$ is either an integer (zero included) or a half-integer. For a given
$s$, the spin component $s_z$ can take the values
$s,s-1,\ldots,-s$, a total of $2s+1$ values. Correspondingly, the wave
function of a particle with spin $s$ has $2s+1$ components.\footnote{For each
species of particle, $s$ is a prescribed number. Consequently, in passing to
the classical limit ($\hbar\to0$), the spin angular momentum $\hbar s$ tends
to zero. This argument is not meaningful for orbital angular momentum,
because $l$ can assume arbitrary values. The classical limit requires
$\hbar$ to tend to zero and $l$ simultaneously to infinity in such a way
that the product $\hbar l$ remains finite.}

\SourcePage{250}{253}
Experiment shows that most elementary particles---electrons, positrons,
protons, neutrons, $\mu$ mesons, and all the hyperons ($\Lambda$, $\Sigma$,
$\Xi$)---have spin $1/2$. There are also elementary particles---the $\pi$
mesons and $K$ mesons---with spin $0$.

The total angular momentum of a particle is the sum of its orbital angular
momentum $\mathbf l$ and its spin $\mathbf s$. Since their operators act on
functions of entirely different variables, they of course commute with each
other.

The eigenvalues of the total angular momentum
\begin{equation}
 \mathbf j=\mathbf l+\mathbf s
 \tag{54.3}
\end{equation}
are determined by the same ``vector-model'' rule as the sum of the orbital
angular momenta of two distinct particles (\S~31). Namely, for prescribed
$l$ and $s$, the total angular momentum can have the values
$l+s,l+s-1,\ldots,|l-s|$. Thus, for an electron (spin $1/2$) with nonzero
orbital angular momentum $l$, the total angular momentum can be
$j=l\pm1/2$; when $l=0$, there is of course only the single value $j=1/2$.

The total-angular-momentum operator $\mathbf J$ of a system of particles is
the sum of the angular-momentum operators $\mathbf j$ for the individual
particles, so that its values are again fixed by the vector-model rules. The
angular momentum $\mathbf J$ can be written
\begin{equation}
 \mathbf J=\sum_a\mathbf j_a=\mathbf L+\mathbf S,
 \qquad \mathbf L=\sum_a\mathbf l_a,\quad
 \mathbf S=\sum_a\mathbf s_a,
 \tag{54.4}
\end{equation}
where $\mathbf S$ may be called the total spin and $\mathbf L$ the total
orbital angular momentum of the system.

Notice that a half-integral (or integral) total spin of a system entails a
half-integral (or integral) total angular momentum, since orbital angular
momentum is always integral. In particular, a system composed of an even
number of identical particles necessarily has integral total spin and hence
integral total angular momentum as well.


The operators of a particle's total angular momentum $\mathbf j$ (or of the
total angular momentum $\mathbf J$ of a system of particles) obey the same
commutation rules as the operators of orbital angular momentum and spin. These
are, in fact, the general commutation rules valid for every angular momentum.
Equation (27.13), obtained from those rules for angular-momentum matrix
elements, is likewise valid for any angular momentum, provided that its matrix
elements are defined with respect to that same angular momentum's
eigenfunctions. With the corresponding changes of notation, formulas
(29.7)--(29.10) for the matrix elements of arbitrary vector quantities also
remain valid.


\SourcePage{251}{254}
\ProblemHeading

A particle of spin $1/2$ is in a state with the definite value $s_z=1/2$.
Find the probabilities of the possible values of the spin projection on an
axis $z'$ inclined at an angle $\theta$ to the $z$ axis.

\SolutionHeading

The mean spin vector $\bar{\mathbf s}$ is evidently directed along the $z$
axis and has magnitude $1/2$. Projecting it onto the $z'$ axis, we find that
the mean spin along $z'$ is $\bar s_{z'}=(1/2)\cos\theta$. On the other hand,
$\bar s_{z'}=(1/2)(w_+-w_-)$, where $w_\pm$ are the probabilities of the
values $s_{z'}=\pm1/2$. Using $w_++w_-=1$ as well, we obtain
\[
 w_+=\cos^2(\theta/2),\qquad w_-=\sin^2(\theta/2).
\]
